\documentclass{article}
\usepackage{amsmath,amssymb,amsthm}
\usepackage[margin=1in]{geometry}

\title{Reading Notes: Closure Certificates\\
HSCC 2024}
\author{Vishnu Murali, Ashutosh Trivedi, Majid Zamani}
\date{Paper Read: 2025-12-28}

\newtheorem{definition}{Definition}
\newtheorem{lemma}{Lemma}
\newtheorem{theorem}{Theorem}
\newtheorem{corollary}{Corollary}

\begin{document}

\maketitle

\section{Paper Metadata}

\begin{itemize}
    \item \textbf{Title}: Closure Certificates
    \item \textbf{Authors}: Vishnu Murali, Ashutosh Trivedi, Majid Zamani
    \item \textbf{Affiliation}: University of Colorado Boulder
    \item \textbf{Publication}: HSCC '24, May 14–16, 2024, Hong Kong
    \item \textbf{DOI}: 10.1145/3641513.3650120
    \item \textbf{arXiv}: 2305.17519v3 [cs.LO] 5 Mar 2024
\end{itemize}

\section{Core Contribution}

This paper introduces \textbf{closure certificates} as a natural extension of barrier certificates from state invariants to transition invariants. Key contributions:

\begin{enumerate}
    \item \textbf{Closure Certificates}: Functional analog of Podelski-Rybalchenko transition invariants
    \item \textbf{Generalization}: Can verify LTL and $\omega$-regular specifications beyond safety
    \item \textbf{SOS and SMT Characterization}: Automate the search for closure certificates
    \item \textbf{Expressiveness}: Show closure certificates can be simpler than barrier certificates for same properties
    \item \textbf{Subsumes State-Triplet}: Prove their approach subsumes Wongpiromsarn et al.'s state-triplet method
\end{enumerate}

\section{System Model}

A discrete-time dynamical system $\mathfrak{S} = (\mathcal{X}, \mathcal{X}_0, f)$ where:
\begin{itemize}
    \item $\mathcal{X} \subseteq \mathbb{R}^n$: state set (compact)
    \item $\mathcal{X}_0 \subseteq \mathcal{X}$: initial states
    \item $f \subseteq \mathcal{X} \times \mathcal{X}$: state transition relation
\end{itemize}

State evolution:
\begin{equation}
\mathfrak{S}: x(t+1) \in f(x(t))
\end{equation}

Labelling function: $\mathcal{L}: \mathcal{X} \rightarrow \Sigma$ maps states to finite alphabet.

\section{Barrier Certificates (Background)}

\begin{definition}[Barrier Certificate]
A function $B: \mathcal{X} \rightarrow \mathbb{R}$ is a barrier certificate for $\mathfrak{S} = (\mathcal{X}, \mathcal{X}_0, f)$ with unsafe states $\mathcal{X}_u$ if:
\begin{align}
B(x) &\leq 0 \quad \forall x \in \mathcal{X}_0 \label{bc1}\\
B(x) &> 0 \quad \forall x \in \mathcal{X}_u \label{bc2}\\
B(x) \leq 0 &\Rightarrow B(x') \leq 0 \quad \forall x \in \mathcal{X}, x' \in f(x) \label{bc3}
\end{align}
\end{definition}

\begin{theorem}[Barrier certificates imply safety]
The existence of a barrier certificate $B$ implies safety (system never reaches $\mathcal{X}_u$).
\end{theorem}

\section{Closure Certificates for Safety}

\begin{definition}[Closure Certificate for Safety]
A function $\mathcal{T}: \mathcal{X} \times \mathcal{X} \rightarrow \mathbb{R}$ is a Closure Certificate (CC) for unsafe states $\mathcal{X}_u$ if $\exists \xi \in \mathbb{R}_{>0}$ such that $\forall x, y \in \mathcal{X}$, $x' \in f(x)$, $x_0 \in \mathcal{X}_0$, $x_u \in \mathcal{X}_u$:
\begin{align}
\mathcal{T}(x, x') &\geq 0 \label{cc1}\\
\mathcal{T}(x', y) \geq 0 &\Rightarrow \mathcal{T}(x, y) \geq 0 \label{cc2}\\
\mathcal{T}(x_0, x_u) &\leq -\xi \label{cc3}
\end{align}
\end{definition}

\textbf{Key Insight}: CC characterizes transition invariant:
\begin{equation}
T = \{(x, x'): \mathcal{T}(x, x') \geq 0 \text{ and } \mathcal{T}(x_0, x) \geq 0 \text{ for some } x_0 \in \mathcal{X}_0\}
\end{equation}

\begin{theorem}[Closure Certificate imply Safety]
The existence of a function $\mathcal{T}: \mathcal{X} \times \mathcal{X} \rightarrow \mathbb{R}$ satisfying (\ref{cc1})-(\ref{cc3}) implies safety.
\end{theorem}

\textbf{Proof idea}: Assume trajectory $\langle x_0, \ldots, x_u, \ldots \rangle$ reaches unsafe state. From (\ref{cc1}): $\mathcal{T}(x_i, x_{i+1}) \geq 0$. From (\ref{cc2}) by induction: $\mathcal{T}(x_0, x_i) \geq 0$ for all $i$. Thus $\mathcal{T}(x_0, x_u) \geq 0$. But (\ref{cc3}) requires $\mathcal{T}(x_0, x_u) \leq -\xi < 0$. Contradiction.

\section{Simplicity of Closure Certificates}

\begin{theorem}[Simplicity of Closure Certificates]
For any natural number $d \in \mathbb{N}$, there exists a system $\mathfrak{S}$ with unsafe set $\mathcal{X}_u$ that cannot be shown safe by a polynomial barrier certificate of degree $d$ but can be shown safe by a linear closure certificate.
\end{theorem}

\textbf{Example} (Figure 1 in paper):
\begin{itemize}
    \item State set: $\mathcal{X} = \{0, 1, 2, 3, 4, 5\}$
    \item Initial states: $\mathcal{X}_0 = \{1, 3, 5\}$
    \item Unsafe states: $\mathcal{X}_u = \{2, 4\}$
    \item Transition: all states go to 0
    \item \textbf{No degree-2 polynomial barrier certificate}: Must change signs at 3+ points (intermediate value theorem)
    \item \textbf{Linear closure certificate exists}: $\mathcal{T}(x, y) = -y$ proves safety
\end{itemize}

\section{Expressiveness}

\begin{theorem}[Expressiveness]
Given a barrier certificate $B: \mathcal{X} \rightarrow \mathbb{R}$, one can compute a closure certificate $\mathcal{T}: \mathcal{X} \times \mathcal{X} \rightarrow \mathbb{R}$.
\end{theorem}

\textbf{Proof}: Let $\gamma \in \mathbb{R}_{>0}$. Define:
\begin{equation}
\mathcal{T}(x, y) = \begin{cases}
0 & \text{if } B(x) > 0 \text{ or } B(y) \leq 0\\
-\gamma & \text{otherwise}
\end{cases}
\end{equation}

This shows closure certificates subsume barrier certificates.

\section{Closure Certificates for Persistence}

\begin{definition}[Closure Certificate for Persistence]
A bounded function $\mathcal{T}: \mathcal{X} \times \mathcal{X} \rightarrow \mathbb{R}$ is a CC for set $\mathcal{X}_{VF}$ (visited finitely often) if $\exists \xi \in \mathbb{R}_{>0}$ such that $\forall x, y \in \mathcal{X}$, $x' \in f(x)$, $x_0 \in \mathcal{X}_0$, $y', y'' \in \mathcal{X}_{VF}$:
\begin{align}
\mathcal{T}(x, x') &\geq 0 \label{ccp1}\\
\mathcal{T}(x', y) \geq 0 &\Rightarrow \mathcal{T}(x, y) \geq 0 \label{ccp2}\\
\mathcal{T}(x_0, y') \geq 0 \land \mathcal{T}(y', y'') \geq 0 &\Rightarrow \mathcal{T}(x_0, y'') \leq \mathcal{T}(x_0, y') - \xi \label{ccp3}
\end{align}
\end{definition}

\textbf{Key mechanism}: Condition (\ref{ccp3}) is a "potential-like" argument - each visit to $\mathcal{X}_{VF}$ decreases potential by $\xi$. Since $\mathcal{T}$ is bounded, only finitely many visits possible.

\begin{theorem}[Closure Certificates imply Persistence]
Existence of $\mathcal{T}$ satisfying (\ref{ccp1})-(\ref{ccp3}) implies traces visit $\mathcal{X}_{VF}$ finitely often.
\end{theorem}

\section{Closure Certificates for LTL Specifications}

Construct product $\mathfrak{S} \otimes \mathcal{A} = (\mathcal{X}', \mathcal{X}_0', f')$ where:
\begin{itemize}
    \item $\mathcal{X}' = \mathcal{X} \times \{0, \ldots, |Q|-1\}$
    \item $\mathcal{X}_0' = \mathcal{X}_0 \times \{q_0 \mid q_0 \in Q_0\}$
    \item $f'((x, q_i)) = \{(x', q_j) \mid q_j \in \delta(q_i, \mathcal{L}(x)), x' \in f(x)\}$
\end{itemize}

\begin{definition}[Closure Certificate for LTL]
A bounded function $\mathcal{T}: \mathcal{X} \times Q \times \mathcal{X} \times Q \rightarrow \mathbb{R}$ is a closure certificate for $\mathfrak{S}$ and NBA $\mathcal{A}$ if $\exists \xi \in \mathbb{R}_{>0}$ such that:
\begin{align}
\mathcal{T}((x, i), (x', i')) &\geq 0 \label{ccltl1}\\
\mathcal{T}((x', i'), (y, j)) \geq 0 &\Rightarrow \mathcal{T}((x, i), (y, j)) \geq 0 \label{ccltl2}
\end{align}
and for all $x_0 \in \mathcal{X}_0$, $s \in Q_0$, $\ell, \ell' \in Acc$:
\begin{equation}
\mathcal{T}((x_0, s), (y, \ell)) \geq 0 \land \mathcal{T}((y, \ell), (y', \ell')) \geq 0 \Rightarrow \mathcal{T}((x_0, s), (y', \ell')) \leq \mathcal{T}((x_0, s), (y, \ell)) - \xi \label{ccltl3}
\end{equation}
\end{definition}

\begin{theorem}[Closure Certificates verify LTL]
Let NBA $\mathcal{A}$ represent $\neg\phi$. Existence of closure certificate satisfying (\ref{ccltl1})-(\ref{ccltl3}) implies $\mathfrak{S} \models_{\mathcal{L}} \phi$.
\end{theorem}

\section{SMT-Based Synthesis}

Use CEGIS (Counterexample-Guided Inductive Synthesis) with template:
\begin{equation}
\mathcal{T}(x, y) = \sum_{m=1}^z c_m p_m(x, y)
\end{equation}

For persistence, sample sets $S_1, S_2 \subseteq \mathcal{X}$, $S_3 \subseteq \mathcal{X}_0$, $S_4 \subseteq \mathcal{X}_{VF}$ and encode:
\begin{align}
\bigwedge_{x \in S_1} \mathcal{T}(x, x') \geq 0\\
\bigwedge_{x \in S_1, y \in S_2} [\mathcal{T}(x', y) \geq 0 \Rightarrow \mathcal{T}(x, y) \geq 0]\\
\bigwedge_{x_0 \in S_3, z, z' \in S_4} [\mathcal{T}(x_0, z) \geq 0 \land \mathcal{T}(z, z') \geq 0 \Rightarrow \mathcal{T}(x_0, z') \leq \mathcal{T}(x_0, z) - \xi]
\end{align}

\textbf{Strengthening}: Can replace implications with:
\begin{align}
\tau_1 \mathcal{T}(x', y) &\leq \mathcal{T}(x, y)\\
\mathcal{T}(x_0, y) - \xi - \mathcal{T}(x_0, y') &\geq \tau_2 \mathcal{T}(x_0, y) + \tau_3 \mathcal{T}(y, y')
\end{align}
for $\tau_1, \tau_2, \tau_3 \in \mathbb{R}_{\geq 0}$. This allows linear programming instead of SMT.

\section{SOS-Based Synthesis}

For semi-algebraic sets $\mathcal{X}, \mathcal{X}_0, \mathcal{X}_{VF}$ defined by polynomials $g_{\mathcal{X}}, g_0, g_{VF}$, construct:
\begin{align}
\sigma_0(x) &= \mathcal{T}(x, x') - \lambda_A^T(x) g_A(x) \label{sos1}\\
\sigma_{inv}(x) &= \mathcal{T}(f(x)) - \beta \mathcal{T}(x) - \lambda_{\mathcal{X}}^T(x) g_{\mathcal{X}}(x) \label{sos2}\\
\sigma_+(x) &= \mathcal{T}(x, y) - \mathcal{T}(f(x), y) - \epsilon - \lambda_+^T(x) g_+(x) \label{sos3}\\
\sigma_-(x) &= \mathcal{T}(x, y) - \mathcal{T}(f(x), y) - \lambda_-^T(x) g_-(x) \label{sos4}
\end{align}

where $\lambda_A, \lambda_{\mathcal{X}}, \lambda_+, \lambda_-$ are SOS multipliers, $\beta \in (0, 1)$, $\epsilon > 0$.

\begin{lemma}
If polynomials $\sigma_0, \sigma_{inv}, \sigma_+, \sigma_-$ are SOS, then $\mathcal{T}$ is a transition persistence certificate.
\end{lemma}

\textbf{Complexity}: $O\left(\binom{2n+d}{d}^2\right)$ per certificate for state dimension $n$ and degree $2d$. For LTL: $O\left(2^{2|\phi|} \times \binom{2n+d}{d}^2\right)$ where $|\phi|$ is formula size.

\section{Subsuming State-Triplet Approach}

\subsection{State-Triplet Method (Wongpiromsarn 2015)}

\textbf{Approach}:
\begin{enumerate}
    \item Consider all simple paths from initial states $Q_0$ to accepting states $Acc$
    \item Break paths into state triplets $(q_m, q_m', q_m'')$ (or edge pairs $(e_m, e_m')$)
    \item Search for barrier certificate to "cut" at least one triplet from each path
    \item If all paths cut, conclude $\mathfrak{S} \models_{\mathcal{L}} \phi$
\end{enumerate}

\textbf{Conservatism}: Treats NBA as finite automaton - cannot handle cases where accepting states are reachable but visited finitely often.

\begin{lemma}
Unrolling NBA more than once does not help find new triplets with different letter pairs.
\end{lemma}

\subsection{Subsumption Result}

\begin{theorem}[CC subsumes State-Triplet]
If barrier certificates $B_1, \ldots, B_k$ show $\mathfrak{S} \models_{\mathcal{L}} \phi$ via state-triplet approach, then there exists a closure certificate $\mathcal{T}$ that also proves $\mathfrak{S} \models_{\mathcal{L}} \phi$.
\end{theorem}

\textbf{Proof sketch}: Divide states into $Q_l$ (left of triplets) and $Q_r$ (right of triplets). Construct $\mathcal{T}$ piecewise:
\begin{itemize}
    \item $\mathcal{T}(x, i, y, j) \geq 0$ if $i \in Q_r$ or both $i, j \in Q_l$
    \item $\mathcal{T}(x, i, y, j) < 0$ if $i \in Q_l$, $j \in Q_r$ (impossible by barrier certificates)
    \item For middle elements of triplets, use corresponding barrier certificate values
\end{itemize}

\section{Case Studies}

\subsection{1D Kuramoto Oscillator (Safety)}

\textbf{System}:
\begin{itemize}
    \item $\mathcal{X} = [0, 2\pi]$
    \item $\mathcal{X}_0 = [\frac{4\pi}{9}, \frac{5\pi}{9}]$
    \item $\mathcal{X}_u = [\frac{7\pi}{9}, \frac{8\pi}{9}]$
    \item $f(x) = x + t_s \Omega + t_s K \sin(-x) - 0.532 x^2 + 1.69$
\end{itemize}

\textbf{Result}: Linear closure certificate $\mathcal{T}(x, y) = 10 - 4.094y$ found in \textbf{10 minutes}. Linear barrier certificate is infeasible.

\subsection{2D Kuramoto Oscillator (LTL)}

Verify $\mathbf{G}\neg a$ using NBA for $\mathbf{F}a$. Found closure certificate in \textbf{1 hour 50 minutes}.

\subsection{Two-Room Temperature (Persistence)}

Verify $a_0 \Rightarrow \mathbf{FG}\neg a_1$ (eventually stay outside comfort zone). Closure certificate found in \textbf{1.5 hours}.

\section{Comparison: Closure Certificates vs. Our Work}

\begin{center}
\begin{tabular}{|p{4cm}|p{5cm}|p{5cm}|}
\hline
\textbf{Aspect} & \textbf{Closure Certificates (Murali 2024)} & \textbf{Our Work} \\
\hline
\textbf{Certificate Type} & $\mathcal{T}: \mathcal{X} \times Q \times \mathcal{X} \times Q \rightarrow \mathbb{R}$ &
\begin{itemize}
\item Safety: $B_k(x,q)$ on $\mathcal{X} \times Q$
\item Persistence: $B_{k,l}(x)$ on $\mathcal{X}$
\end{itemize} \\
\hline
\textbf{Search Space Dimension} & $\mathcal{X} \times Q \times \mathcal{X} \times Q$ (4D: two state copies, two automaton states) &
\begin{itemize}
\item Safety: $\mathcal{X} \times Q$ (2D)
\item Persistence: $\mathcal{X}$ (1D)
\item Total: at most $\mathcal{X} \times Q \times Q$ (3D)
\end{itemize} \\
\hline
\textbf{Number of Certificates} & 1 piecewise certificate with $|Q|^2$ components &
\begin{itemize}
\item $|Acc^{start}|$ safety certificates
\item OR $|Acc^{pair}|$ persistence certificates per start state
\end{itemize} \\
\hline
\textbf{Verification Approach} & Prove accepting states visited finitely often via ranking on product &
\begin{itemize}
\item Safety: Block reachability to $q_k \in Acc^{start}$
\item Persistence: Ranking function on $\mathcal{X}$ alone
\end{itemize} \\
\hline
\textbf{Synthesis Time} &
\begin{itemize}
\item 1D Kuramoto: 10 min
\item 2D Kuramoto: 110 min
\item Two-room: 90 min
\end{itemize} &
\begin{itemize}
\item 1D Kuramoto: 4.25 s
\item 2D Kuramoto: 7.27 s
\item Two-room: 1.27 s (SOS), 8.77 s (SMT)
\end{itemize} \\
\hline
\textbf{Speedup} & Baseline & \textbf{141× to 4248×} \\
\hline
\textbf{Accepting Conditions} & State-based NBA & Transition-based NBA \\
\hline
\textbf{Conservatism} & Lower than state-triplet (allows finite visits) & Even lower (separate safety/persistence) \\
\hline
\textbf{Completeness} & Incomplete (sound only) & Incomplete (sound only) \\
\hline
\end{tabular}
\end{center}

\section{Key Insights for Our Paper}

\subsection{Our Advantages Over Closure Certificates}

\begin{enumerate}
    \item \textbf{Much Smaller Search Space}:
    \begin{itemize}
        \item Closure: $\mathcal{X} \times Q \times \mathcal{X} \times Q$ (cartesian product of state pairs with automaton state pairs)
        \item Our safety: $\mathcal{X} \times Q$ per certificate
        \item Our persistence: $\mathcal{X}$ per certificate
        \item Total complexity: $\mathcal{X} \times Q \times Q$ vs. $\mathcal{X} \times \mathcal{X} \times Q \times Q$
    \end{itemize}

    \item \textbf{Dramatic Speedup}: 141× to 4248× faster synthesis on same benchmarks

    \item \textbf{Modular Certificates}:
    \begin{itemize}
        \item Closure: 1 piecewise certificate with $|Q|^2$ components
        \item Our: separate certificates for each accepting transition, easier to synthesize
    \end{itemize}

    \item \textbf{Persistence on State Space Alone}: Our persistence certificates $B_{k,l}(x)$ only depend on $x \in \mathcal{X}$, not on automaton states

    \item \textbf{Transition-Based NBA}: More efficient for LTL verification than state-based
\end{enumerate}

\subsection{Positioning in Related Work}

When citing this paper in our related work:

\begin{quote}
``Murali et al.~\cite{murali2024closure} introduced closure certificates as functional analogs of transition invariants, generalizing barrier certificates to verify LTL specifications. Their closure certificates are defined over $\mathcal{X} \times Q \times \mathcal{X} \times Q$, creating high-dimensional search spaces with heavy computational burdens. Our certificates operate on smaller spaces: persistence certificates on $\mathcal{X}$, safety certificates on $\mathcal{X} \times Q$, yielding a total search space of at most $\mathcal{X} \times Q \times Q$—significantly smaller than $\mathcal{X} \times Q \times \mathcal{X} \times Q$. Case studies show that for benchmarks requiring closure certificates, our certificates synthesize in seconds, achieving 141× to 4248× speedup.''
\end{quote}

\subsection{Example Demonstrating Advantage}

For the 2D Kuramoto example:
\begin{itemize}
    \item \textbf{Closure certificate method}: 110 minutes (6600 seconds)
    \item \textbf{Our method}: 7.27 seconds
    \item \textbf{Speedup}: 909×
    \item \textbf{Reason}: Search space $\mathcal{X} \times Q$ (2D state × 2 automaton states = 4D) vs. $\mathcal{X} \times Q \times \mathcal{X} \times Q$ (2D × 2 × 2D × 2 = 8D)
\end{itemize}

\section{Citation Correctness}

\textbf{Verified}: This paper is \textbf{directly relevant and essential} to our work:
\begin{itemize}
    \item Uses certificates for LTL verification ✓
    \item Works with NBA ✓
    \item Extends barrier certificates to temporal logic ✓
    \item Is our main comparison baseline ✓
    \item Provides benchmark case studies we compare against ✓
\end{itemize}

\textbf{Recommendation}: This citation is \textbf{absolutely critical}. It should be cited as the main competing approach and our primary baseline for experimental comparison. All speedup numbers in our paper reference this work.

\section{Summary}

Closure certificates bridge barrier certificates and transition invariants, enabling LTL verification through well-founded arguments on state pairs. While less conservative than state-triplet methods, they operate on high-dimensional spaces ($\mathcal{X} \times Q \times \mathcal{X} \times Q$), leading to significant computational overhead. Our work achieves the same verification goals with drastically smaller search spaces and 2-3 orders of magnitude speedup.

\end{document}
